\documentclass{article}
\usepackage{graphicx} % Required for inserting images
\usepackage{fancyhdr}
\usepackage{amsmath}
\usepackage{float}


\begin{document}
\pagestyle{fancy}
\fancyhead[RO]{Mehmet Çağatay Kurmaç}
\fancyhead[LO]{Static VB Example}


\section*{1.}
Let's begin by the intention behind EM algorithm (Will be usefull for 3rd question.)

For any measurement $y$ we can write the likelihood function with respect to the parameter $\lambda$ as follows,where $x$ is the latent variable:
\[
\begin{aligned} 
p(y|\lambda) &= \int p(y,x|\lambda) dx = \int p(y,x|\lambda) *\frac{p(x|y,\lambda)}{p(x|y,\lambda)} dx \\
            &= E_{x|y,\lambda}\left[ \frac{p(y,x|\lambda)}{p(x|y,\lambda)} \right] \\
\end{aligned}
\]

\subsubsection*{Jensen's Inequality}   
\[
  ln E[X] \geq E[ln X]
\]

According to Jensen's inequality, we can write the following:
\[
\begin{aligned}
    ln p(y|\lambda) &= ln E_{x|y,\lambda}\left[ \frac{p(y,x|\lambda)}{p(x|y,\lambda)} \right] \\
    &\geq E_{x|y,\lambda}\left[ ln \frac{p(y,x|\lambda)}{p(x|y,\lambda)} \right] \\
\end{aligned}
\]
We denote the $E_{x|y,\lambda}\left[ ln \frac{p(y,x|\lambda)}{p(x|y,\lambda)} \right]$ as Evidence Lower Bound. We can factor it further thanks to natural logarithm.

\begin{equation}
    ln p(y|\lambda) \geq E_{x|y,\lambda}\left[ ln p(y,x|\lambda)\right] - E_{x|y,\lambda}\left[ ln p(x|y,\lambda) \right]
\end{equation}

First element is our Q function, and the second element is the entropy of the posterior distribution. Since the entropy of the posterior distribution
does not depend on new values of $\lambda$, we will regard it as a constant later on.

Our idea is to approximate the distribution of the latent variable $x$ with respect to the current value of parameters and latent variables. (E - step)
Then, use this distribution to maximize the Q function with respect to the parameters and latent variables. (M - step)

In our expectation step, we will come up with an expression for the posterior distribution $p(x|y,\lambda)$. Calculate its value with respect
to the current value of $\lambda$. Then, in the maximization step, we will maximize the Q function featured here, with respect to parameters and 
latent variables.

For unknown parameters and systems with unobserved latent variables, This allows us to calculate the maximum likelihood estimates of the parameters iteratively,
without needing to compute the posterior distribution of the latent variables, which is often intractable.

In our problem, we can estimate the lambda with the following steps:\\
\textbf{E - Step}:\\
Notice that our variables $x$ and $v$ are gaussian random variables. Which makes our system linear and Gaussian. Therefore, we can use Kalman
filter to estimate the posterior distribution of the latent variable $x$ given the measurement $y$, given the current estimate or the initial
value of the parameter $\lambda$. 
\[
    \begin{aligned}
        p(x|y,\lambda) &= \frac{p(y|x,\lambda)p(x)}{p(y|\lambda)} \\
                    K &= \frac{\sigma_x^2}{\sigma_x^2 + \frac{1}{\lambda}}\\
                     \mu_{x|y,\lambda} &= \bar{x} + K(y - \bar{x}) = \frac{\lambda*\sigma_x^2*y+\bar{x}}{\lambda*\sigma_x^2+1} \\
                     \sigma_{x|y,\lambda} &= \sigma_x^2 - K\sigma_x^2 = \frac{\sigma_x^2}{\lambda*\sigma_x^2+1}
    \end{aligned}
\]

Thus: 

\[
    p(x|y,\lambda) = \mathcal{N}\left(x;\frac{\lambda*\sigma_x^2*y+\bar{x}}{\lambda*\sigma_x^2+1}, \frac{\sigma_x^2}{\lambda*\sigma_x^2+1}\right)
\]
\\
\textbf{M - Step}:\\

Resulting Q function is as follows:
\[
\begin{aligned}
    E_{x|y,\lambda}\left[ ln p(y,x|\lambda)\right] &= E_{y|x,\lambda}\left[ ln p(y|x,\lambda)\right] + E_{x|y,\lambda}\left[ ln p(x)\right] \\
    p(y|x,\lambda) &= \mathcal{N}(y;x,\frac{1}{\lambda}) = \sqrt{\frac{\lambda}{2\pi}} e^{-\frac{\lambda}{2}(y-x)^2} \\
    ln p(y|x,\lambda) &= \frac{1}{2} ln \lambda - \frac{1}{2} ln 2\pi - \frac{\lambda}{2}(y-x)^2 \\
    E_{x|y,\lambda}\left[ ln p(y|x,\lambda)\right] &= \frac{1}{2} \ln \lambda - \frac{1}{2} \ln 2\pi - \frac{\lambda}{2} E_{x|y,\lambda}\left[ (y-x)^2 \right]
\end{aligned}
\]
Regarding the terms independent of $\lambda$ as constant (Entire $p(x)$ expectation is also independent of $\lambda$), since we will take the derivative with respect to $\lambda$, we can write the following:
\[
\begin{aligned}
    E_{x|y,\lambda}\left[ ln p(y,x|\lambda)\right] &= \frac{1}{2} \ln \lambda - \frac{\lambda}{2} E_{x|y,\lambda}\left[ (y-x)^2 \right] + C \\
                                                   &= \frac{1}{2} \ln \lambda - \frac{\lambda}{2} \left( (y-\mu_{x|y,\lambda})^2 + \sigma_{x|y,\lambda}^2 \right) + C
\end{aligned}
\]
Taking derivative with respect to lambda and setting it to zero, we can find the optimal value of $\lambda$:
\[
\begin{aligned}
    \frac{d}{d\lambda} E_{x|y,\lambda}\left[ ln p(y,x|\lambda)\right] &= \frac{1}{2\lambda} - \frac{1}{2} \left( (y-\mu_{x|y,\lambda})^2 + \sigma_{x|y,\lambda}^2 \right) = 0 \\
    \lambda &= \frac{1}{(y-\mu_{x|y,\lambda})^2 + \sigma_{x|y,\lambda}^2} 
\end{aligned}
\]

Following from here, we can iteratively update the value of $\lambda$ until convergence.
\section*{2.}

Let's also begin with the formulation of the Variational Bayes framework. 

\[
    \begin{aligned}
        p(y,x|\lambda) &= p(x|y,\lambda)*p(y|\lambda) \\
        ln(p(y,x|\lambda)) &= ln(p(x|y,\lambda)+ ln(p(y|\lambda))) \\
        ln(p(y|\lambda)) &= ln(p(y,x|\lambda)) - ln(p(x|y,\lambda)) \\
    \end{aligned}
\]
Now, we will add and substract an arbitrary distribution $q(x)$ to the right hand side of the equation:
\[
    \begin{aligned}
        ln(p(y|\lambda)) &= ln(p(y,x|\lambda)) - ln(p(x|y,\lambda)) + ln(q(x)) - ln(q(x))\\
                         &= ln(\frac{p(y,x|\lambda)}{q(x)}) - ln(\frac{p(x|y,\lambda)}{q(x)}) \\
    \end{aligned}
\]
Take the expectation of both sides with respect to $q(x)$:
\begin{equation}
    ln(p(y|\lambda)) = E_{q(x)}\left[ ln(\frac{p(y,x|\lambda)}{q(x)}) \right] - E_{q(x)}\left[ ln(\frac{p(x|y,\lambda)}{q(x)}) \right]
\end{equation}

The first term expresses the Evidence Lower Bound (ELBO) and the second term is the Kullback-Leibler divergence between the arbitrary distribution
$q(x)$ and the posterior distribution $p(x|y,\lambda)$. Since KL divergence is always non-negative, we can write the following:
\[
    ln(p(y|\lambda)) \geq E_{q(x)}\left[ ln(\frac{p(y,x|\lambda)}{q(x)}) \right]
\] 

By maximizing the ELBO with respect to $q(x)$, we can find the best approximation of the posterior distribution $p(x|y,\lambda)$.
If the arbitrary distribution $q(x)$ is a factorized distribution $q(x) = \prod_{i=1}^{n} q_i(x_i)$, We can formulate the ELBO as follows:

\[
    \begin{aligned}
        \mathcal{L}(q) &= \sum_{X}\prod_{i=1}^{n} q_i(x_i) \left( \ln p(y,x|\lambda) - \sum_{i=1}^{n} \ln q_i(x_i) \right) \\
                       &= \sum_{x_j}q_j(x_j) \left(\sum_{X_{i\neq j}}\prod_{i\neq j} q_i(x_i)*\ln p(y,x|\lambda) - \sum_{X_{i\neq j}}\prod_{i\neq j} q_i(x_i)\sum_{k\neq j} \ln q_k(x_k) \right) - \sum_{x_j}q_j(x_j)\ln q_j(x_j) \\
    \end{aligned}
\]
For $k\neq i$, $\sum_{X_{i}}q_i(x_i)\ln q_k(x_k) = \ln q_k(x_k)$, since $\sum_{X_{i}}q_i(x_i) = 1$. Thus, it acts as a constant with respect to
component $q_j(x_j)$. Also, we can express the first term as $E_{i\neq j}\left[ \ln p(y,x|\lambda) \right]$. Therefore, we can write the following:
\[
    \begin{aligned}
        \mathcal{L}(q) &= \sum_{x_j}q_j(x_j) \left( E_{i\neq j}\left[ \ln p(y,x|\lambda) \right]\right) - \sum_{x_j}q_j(x_j)\ln q_j(x_j)  \\
    \end{aligned}
\]  
If we denote the first term as $E_{i\neq j}\left[ \ln p(y,x|\lambda) \right] = ln(\tilde{p}(x_j))$, we can write the following:
\[
    \begin{aligned}
        \mathcal{L}(q) &= \sum_{x_j}q_j(x_j) \left( ln(\tilde{p}(x_j))\right) - \sum_{x_j}q_j(x_j)\ln q_j(x_j)  \\
                       &= -\sum_{x_j}q_j(x_j) \left( \ln(\frac{q_j(x_j)}{\tilde{p}(x_j)})\right) \\
    \end{aligned}
\]
It basically became a negative KL divergence between $q_j(x_j)$ and $\tilde{p}(x_j)$. Since KL divergence is always non-negative, 
we can maximize the ELBO by minimizing the KL divergence. KL divergence is minimized to zero when $q_j(x_j) = \tilde{p}(x_j)$. Therefore
, we can find optimal $q_j(x_j)$ as follows:
\[
    q_j^*(x_j) = \tilde{p}(x_j) = exp(E_{i\neq j}\left[ \ln p(y,x|\lambda) \right])
\]

Let's apply this to our problem. Distribution of $\lambda$ is given as $\lambda \sim Gam(\lambda|a,b)$. The factorization of the latent variable and parameters is 
given as $q(x,\lambda) = q_x(x)q_\lambda(\lambda)$. Therefore, we can write the following:

\[
    \begin{aligned}
        ln(q_x^*(x)) &= E_{\lambda}\left[ \ln p(y,x,\lambda) \right] \\
                     &= E_{\lambda}\left[ \ln p(y|x,\lambda) + \ln p(x) + \ln p(\lambda) \right] \\
                        &= E_{\lambda}\left[ \ln p(y|x,\lambda) \right] + \ln p(x) + E_{\lambda}\left[ \ln p(\lambda) \right] \\
                        &= E_{\lambda}\left[ \ln p(y|x,\lambda) \right] + \ln p(x) + C \\
                        &= E_{\lambda}\left[ \ln \mathcal{N}(y;x,\frac{1}{\lambda}) \right] \\
                        &= E_{\lambda}\left[ \frac{1}{2} ln \lambda - \frac{1}{2} ln 2\pi - \frac{\lambda}{2}(y-x)^2 \right] + \ln p(x) + C \\
                        &= \frac{1}{2} E_{\lambda}\left[ ln \lambda \right] - \frac{1}{2} E_{\lambda}\left[ \lambda \right](y-x)^2 + \ln p(x) + C \\
                        &= - \frac{a}{2b}(y-x)^2 - \frac{1}{2} ln 2\pi - \frac{(x-\bar{x})^2}{2\sigma_x^2} + C \\
                        &= -\frac{a}{2b}(y-x)^2 - \frac{(x-\bar{x})^2}{2\sigma_x^2} \\
                        &= -\frac{1}{2}\left( \frac{a}{b}(y-x)^2 + \frac{(x-\bar{x})^2}{\sigma_x^2} \right) \\
                        &= -\frac{1}{2}\left( \left(\frac{a}{b} + \frac{1}{\sigma_x^2}\right)x^2 - 2\left(\frac{a}{b}y + \frac{\bar{x}}{\sigma_x^2}\right)x + \left(\frac{a}{b}y^2 + \frac{\bar{x}^2}{\sigma_x^2}\right) \right) \\
                        &= -\frac{1}{2}\left( \left(\frac{a}{b} + \frac{1}{\sigma_x^2}\right)x^2 - 2\left(\frac{a}{b}y + \frac{\bar{x}}{\sigma_x^2}\right)x \right) + C \\
    \end{aligned}
\]

The last expression is similar to the natural logarithm of a Gaussian distribution, which is $\frac{1}{2}ln(\lambda)-\frac{1}{2}\ln(2\pi) - \frac{\lambda}{2}(x-\mu)^2$. If we
remove the constant terms, it becomes $- \frac{\lambda}{2}(x-\mu)^2$, which has the form $- \frac{\lambda}{2}x^2 + \lambda\mu x + C$. By comparing the coefficients, 
we can find the mean and variance of optimal distribution of $x$ as follows:
\[
    \begin{aligned}
        \sigma_x^2 &= \frac{1}{\frac{a}{b} + \frac{1}{\sigma_x^2}} \\
                 &= \frac{b\sigma_x^2}{a\sigma_x^2 + b} \\
        \mu_x &= \frac{\frac{a}{b}y + \frac{\bar{x}}{\sigma_x^2}}{\frac{a}{b} + \frac{1}{\sigma_x^2}} \\
    \end{aligned}
\]

Let's also find the optimal distribution of $\lambda$:

\[
    \begin{aligned}
        ln(q_\lambda^*(\lambda)) &= E_{x}\left[ \ln p(y,x,\lambda) \right] \\
                     &= E_{x}\left[ \ln p(y|x,\lambda) + \ln p(x) + \ln p(\lambda) \right] \\
                        &= E_{x}\left[ \ln p(y|x,\lambda) \right] + E_{x}\left[ \ln p(x) \right] + \ln p(\lambda) \\
                        &= E_{x}\left[ \ln p(y|x,\lambda) \right] + C + ln(p(\lambda))\\
                        &= E_{x}\left[ ln(\mathcal{N}(y;x,\frac{1}{\lambda}))\right] + ln(Gam(\lambda|a,b)) + C\\
                        &= E_{x}\left[ \frac{1}{2} ln \lambda - \frac{1}{2} ln 2\pi - \frac{\lambda}{2}(y-x)^2\right] - ln(\Gamma(a)) + a ln(b) + (a-1)ln(\lambda) - b\lambda + C\\
                        &= \frac{1}{2} ln \lambda - \frac{\lambda}{2} E_{x}\left[ (y-x)^2\right] + (a-1)ln(\lambda) - b\lambda + C\\
                        &= \left(\frac{1}{2} + a - 1\right)ln(\lambda) - \left(b + \frac{1}{2} E_{x}\left[ (y-x)^2\right]\right)\lambda + C \\
                        &= \left(a - \frac{1}{2}\right)ln(\lambda) - \left(b + \frac{1}{2} E_{x}\left[ (y-x)^2\right]\right)\lambda + C \\
    \end{aligned}
\]

The resulting expression is similar to the natural logarithm of a Gamma distribution, which is $(a-1)ln(\lambda) - b\lambda - a ln(b) - ln(\Gamma(a))$. By comparing the coefficients
, we can find the shape and scale parameters of optimal distribution of $\lambda$ as follows:
\[
    \begin{aligned}
        a^* &= a + \frac{1}{2} \\
        b^* &= b + \frac{1}{2} E_{x}\left[ (y-x)^2\right] \\
    \end{aligned}
\]  

Then the estimated $\lambda$ is basically the mean of the optimal distribution of $\lambda$, which is $\frac{a^*}{b^*}$.

\section*{3.}

As we can see from equations (1) and (2) (Only numbered equations), EM and VB are equivalent, when our latent variable distribution is 
equal to the posterior distribution of the latent variable. 

\[ 
    q(x) = p(x|y,\lambda);
\]

This makes KL divergence term zero and ELBO term exactly matches the EM algorithm's Q function + entropy of the posterior distribution.

\end{document}